\begin{figure}[ht]
  \centering
  \begin{tikzpicture}
    \begin{axis}[paperaxis,
      xlabel={Année}, ylabel={Écart France moins comparables (levier)},
      xtick={2013,2014,2015,2016,2017,2018,2019,2020,2021,2022},
      xticklabel style={/pgf/number format/1000 sep=, rotate=45, anchor=east},
    ]
      \draw[zeroline, line width=0.6pt] (axis cs:2012.5,0) -- (axis cs:2022.5,0);
      % La France, en bleu, avec intervalle a 95 pour cent.
      \addplot[franceblue, mark=*, mark options={scale=0.8}, error bars/.cd,
        y dir=both, y explicit, error bar style={franceblue}]
        table[col sep=comma, x=year, y=coef, y error expr=1.96*\thisrow{std_error}]
        {figures/levier_event.csv};
      % Reforme de 2018 : marqueur vert, la couleur de la baisse d'impot.
      \draw[baisse2018, dashed, line width=1pt]
        ({axis cs:2017.5,0}|-{rel axis cs:0,0}) -- ({axis cs:2017.5,0}|-{rel axis cs:0,1});
    \end{axis}
  \end{tikzpicture}
  \caption{Étude d'événement du levier autour de la réforme de 2018}
  \label{fig:levier}
  \fignote{Coefficients $\tau_j$ de l'équation~\eqref{eq:event} pour le levier, total des dettes sur
  total de l'actif, écart de la France aux cinq pays comparables l'année $j$, année de référence 2017 omise.
  Barres d'erreur à quatre-vingt-quinze pour cent sur l'erreur-type regroupée par pays, indicative.
  La ligne verte verticale marque la réforme de 2018. Le levier français, un peu supérieur en début
  de période, reflue vers celui des pays comparables et ne remonte pas après la réforme. Ce levier
  agrège l'ensemble des dettes ; la hausse de la dette financière garantie pendant la crise
  sanitaire, logée en trésorerie, n'y apparaît pas et est caractérisée à la
  section~\ref{sec:robustesse}.}
\end{figure}
